\section{The conormal rigidity theorem}\label{sec:general-tomography}
Fix the endpoint-free native experiment of Section~\ref{sec:binary-hankel}, with $k$ interior double roots, $2k+1$ clocks, and unknown positive exposures. Fix $d\ge2$ and nonzero loadings $a_i\in\R^d$ such that
\[
 \mathcal A_A:\Sym_d\longrightarrow\R^k,\qquad
 S\longmapsto(a_i^{\mathsf T}Sa_i)_i
\]
is injective. Put $q_A(\eta)=\mathcal A_A(\eta\eta^{\mathsf T})$ and
$\mathcal V_A=q_A(\R^d)$. For $R>0$ let
\[
 C_R(b)=\min_{\eta\in\R^d}(q_A(\eta)-b)^{\mathsf T}R^{-1}(q_A(\eta)-b),
 \qquad b\in\R_+^k.
\]
The full probability residual also includes the free score directions, which have been profiled out in this definition. Theorem~\ref{thm:synchronized} supplies its statistical interpretation. We call $\eta$ admissible when $\eta\ne0$ and every $a_i^{\mathsf T}\eta\ne0$; then $q_A(\eta)$ is a strictly positive target.

At such a point set
\[
 N_\eta=\ker Dq_A(\eta)^{\mathsf T},\qquad
 H_R(\eta)=\tfrac12D^2 C_R(q_A(\eta)).
\]
Let $\Pi_\eta$ be the Euclidean orthogonal projection onto $N_\eta$. The quadratic form $H_R(\eta)$ restricted to $N_\eta$ is positive definite. For $n\in N_\eta$, its \emph{dual second-jet value} is
\begin{equation}\label{eq:dual-jet}
 h_R(\eta,n)=n^{\mathsf T}
       \big(H_R(\eta)|_{N_\eta}\big)^{-1}n.
\end{equation}
This is determined by the observed second jet and known tangent space. In codimension greater than one it is not, in general, the reciprocal of one directional second derivative. The finite query count below is for the scalar oracle \eqref{eq:dual-jet}; the three-root case in Section~\ref{sec:tomography} reduces it to ordinary directional curvatures.

Write $\mathscr R_k$ for the native inverse-information cone and $\mathscr W_k$ for its linear span. If $d=2$, define the symmetric matrix $B_A$ by
\[
 (B_A)_{ii}=0,\qquad (B_A)_{ij}=\det(a_i,a_j)^2\quad(i\ne j),
\]
and, for $i<j<l$, put
\begin{equation}\label{eq:triple-transversality}
 \tau_{ijl}(A)=\frac{(r_j-r_i)(B_A)_{ij}}{d_i d_j}
 -\frac{(r_l-r_i)(B_A)_{il}}{d_i d_l}
 +\frac{(r_l-r_j)(B_A)_{jl}}{d_j d_l}.
\end{equation}
The $d_i$ here are the interpolation coefficients of \eqref{eq:cauchy-columns}, not the dimension $d$.

\begin{theorem}[Conormal rigidity in native Hankel geometry]\label{thm:general-tomography}
All loadings, sites, clocks, and central channel parameters are fixed as above.

\smallskip\noindent\textup{(a) Complete second-jet ambiguity.}
If $d=2$, two positive definite matrices $R,R'$ give identical second contact jets at every admissible point if and only if $R'-R\in\R B_A$. If $d\ge3$, identical second contact jets imply $R'=R$. These are statements for arbitrary positive metrics, before imposing native information geometry.

\smallskip\noindent\textup{(b) Exact native criterion.}
The native image has dimension $2k-1$, and
\begin{equation}\label{eq:all-hankel-relations}
 \mathscr W_k=\left\{X\in\Sym_k:
 \frac{(r_j-r_i)X_{ij}}{d_i d_j}
 -\frac{(r_l-r_i)X_{il}}{d_i d_l}
 +\frac{(r_l-r_j)X_{jl}}{d_j d_l}=0\ (i<j<l)\right\}.
\end{equation}
For $d=2$, the full second-jet observation is injective on $\mathscr R_k$ exactly when at least one $\tau_{ijl}(A)$ is nonzero. When all vanish its fibres are exactly affine lines parallel to $B_A$ intersected with $\mathscr R_k$. For $d\ge3$ it is always injective on $\mathscr R_k$.

\smallskip\noindent\textup{(c) Finite recovery and contact-data fibres.}
In each injective case, $2k-1$ suitably chosen scalar dual second-jet values recover $R$ linearly, and no smaller number suffices, even for deterministic adaptive choices. They therefore determine $C_R(b)$ for every $b\ge0$. At unit total exposure the attainable vectors of these $2k-1$ dual values are a linear image of $\{y\in\mathscr M:\chi(y)\le1\}$. Their relative boundary has a positive-Hessian analytic defining function. Each boundary datum has one positive-exposure design; each interior datum has a design fibre diffeomorphic to $S^1$. The same sharp query lower bound holds on the budget interior.

Recovery is Lipschitz with the reciprocal of the smallest singular value of the selected linear measurement operator on $\mathscr W_k$. On compact positive-metric families it controls $C_R$ uniformly on compact target sets. For $d=2$, transversality in \eqref{eq:triple-transversality}, and not positive definiteness alone, removes the otherwise unavoidable one-dimensional ambiguity.
\end{theorem}

\begin{lemma}[Normal restriction of inverse information]\label{lem:normal-restriction}
At every admissible point,
\begin{equation}\label{eq:normal-restriction}
 H_R(\eta)=\Pi_\eta
   \big(\Pi_\eta R\Pi_\eta|_{N_\eta}\big)^{-1}\Pi_\eta,
 \qquad h_R(\eta,n)=n^{\mathsf T}Rn.
\end{equation}
\end{lemma}
\begin{proof}
Injectivity of $\mathcal A_A$ identifies the image of rank-one positive matrices away from zero with a smooth embedded manifold of dimension $d$. The map is proper, since $\|q_A(\eta)\|\ge c|\eta|^2$. Squared distance in the fixed metric is smooth in a local tubular neighbourhood. To verify that the local projection suffices globally, the trial at the base point bounds the minimum by the squared displacement. Positive definiteness then forces every minimizing image to approach that base point; properness rules out escape. In local coordinates, the stationary equations have the positive tangent Gram matrix as their derivative, so the implicit function theorem gives a unique smooth minimizer.

The quadratic term is obtained by minimizing over the tangent space. If $u$ is a displacement, the minimizing residual $w$ satisfies $R^{-1}w\in N_\eta$ and $\Pi_\eta w=\Pi_\eta u$. Writing $w=Rn$ gives
$\Pi_\eta R\Pi_\eta n=\Pi_\eta u$ and the displayed quadratic form. Inverting its normal restriction proves the dual identity.
\end{proof}

\begin{lemma}[Quadratic forms on the conormal union]\label{lem:determinantal-kernel}
Let $\mathcal A_A^*(n)=\sum_i n_i a_i a_i^{\mathsf T}$. A quadratic form $F(n)$ on $\R^k$ vanishes on $N_\eta$ for every admissible $\eta$ if and only if it is zero when $d\ge3$, or is a multiple of $\det\mathcal A_A^*(n)$ when $d=2$.
\end{lemma}
\begin{proof}
The adjoint is onto $\Sym_d$. Moreover $n\in N_\eta$ is equivalent to $\mathcal A_A^*(n)\eta=0$. Consequently the union of these normal spaces is dense in the inverse image of the singular symmetric matrices. For clarity, if a singular matrix $S$ has a null vector on an excluded loading hyperplane, choose orthogonal maps $U_j\to I$ taking that vector off all finitely many excluded hyperplanes. Then $U_jSU_j^{\mathsf T}$ has an admissible null vector. A fixed linear right inverse of $\mathcal A_A^*$ lifts these perturbations and preserves convergence to any prescribed lift of $S$. Thus continuity extends vanishing to all lifts of singular matrices.

Choose a linear splitting $n=L(S)+z$, $z\in\ker\mathcal A_A^*$. The term purely quadratic in $z$ vanishes by taking $S=0$. The mixed term vanishes by taking singular $S$ and varying $z$: singular symmetric matrices span $\Sym_d$ for $d\ge2$, since rank-one matrices do. It remains to classify quadratic forms $f(S)$ vanishing on singular symmetric matrices.

For $d\ge3$, fix an orthogonal matrix $U$ and restrict to $S=U\diag(t_1,\ldots,t_d)U^{\mathsf T}$. The resulting degree-two polynomial vanishes whenever any $t_i=0$. Every monomial has degree at most two and hence omits some variable when $d\ge3$; its coefficient must therefore vanish. Every symmetric matrix has such an orthogonal diagonalization, so $f=0$.

For $d=2$, write $S=\left(\begin{smallmatrix}x&y\\y&z\end{smallmatrix}\right)$. Substituting $S=(u,v)(u,v)^{\mathsf T}$ into a general quadratic form shows that the coefficients of $x^2,z^2,xy,yz$ vanish and that the coefficients of $xz,y^2$ are opposite. Thus $f$ is a multiple of $xz-y^2$. Conversely this determinant vanishes on the entire normal union. Finally
\[
 2\det\mathcal A_A^*(n)
   =2\sum_{i<j}\det(a_i,a_j)^2n_i n_j=n^{\mathsf T}B_A n.
\]
This identifies the nonzero ambiguity direction when $d=2$.
\end{proof}

\begin{proof}[Proof of Theorem~\ref{thm:general-tomography}]
By Lemma~\ref{lem:normal-restriction}, equality of second jets is equivalent to equality of the restrictions of $R$ to all normal spaces. Polarization and Lemma~\ref{lem:determinantal-kernel} give (a).

For each clock, the Cauchy column $c_i(T)=d_i/(T-r_i)$ satisfies every three-site relation in \eqref{eq:all-hankel-relations}, by multiplying through the three denominators. Positive sums therefore satisfy them. Conversely those linear relations specify a space of dimension $2k-1$: the $k$ diagonal entries are arbitrary, and the off-diagonal entries are specified by $k-1$ numbers $u_2,\ldots,u_k$ through
\[
 \frac{X_{ij}}{d_i d_j}=\frac{u_j-u_i}{r_j-r_i},\qquad u_1=0.
\]
The native span has the same dimension by Theorem~\ref{thm:binary-hankel}, proving (b). In dimension two, its intersection with the ambiguity line is zero precisely when at least one displayed three-site functional is nonzero on $B_A$. Relative openness of the native cone gives the asserted nontrivial line fibres in the exceptional case.

In every injective case the scalar functionals $X\mapsto n^{\mathsf T}Xn$, for accessible conormals $n$, separate $\mathscr W_k$. A finite basis of $2k-1$ such functionals therefore exists. Their linear inversion proves finite recovery. A deterministic adaptive procedure using fewer queries leaves a nonzero common kernel at any interior reference metric. Small opposite perturbations along it give the same transcript and remain native; this proves the lower bound.

The moment map has corank $(2k+1)-(2k-1)=2$. Lemma~\ref{lem:reciprocal-general} and Theorem~\ref{thm:budget} give the budget image and its singleton or spherical fibres. The selected linear measurements are an isomorphism from moment space, so they transport that image, its positive defining Hessian, and its design fibres exactly. The budget interior is relatively open in $\mathscr W_k$, giving the same lower bound. Finally finite-dimensional inversion gives the singular-value estimate. Comparing minimizing residuals in two uniformly positive metrics gives uniform control of the entire minimum, without requiring unique minimizers; the explicit estimate is supplied in the proof of Corollary~\ref{thm:tomography} below.
\end{proof}
